\chapter{Delen en veelvouden}\label{ch:g5:division}

De staartdeling uit \cref{ch:g4:division} krijgt twee uitbreidingen: \omterm{def:g5:division:multiple}{delers}
van twee cijfers, en \omterm{def:g3:division:remainder}{quotiënten} die voorbij het \omterm{def:g4:decimals:point}{decimaalteken} doorgaan ---
$3 \div 4$ krijgt eindelijk een antwoord, $0.75$. Het hoofdstuk eindigt met
\omterm{def:g5:division:multiple}{veelvouden}, \omterm{def:g5:division:multiple}{delers} en de eerste deelbaarheidsregels.

\section{Delen door een getal van twee cijfers}

\begin{method}[Delers van twee cijfers]\label{met:g5:division:twodigits}
Dezelfde manier als eerder, met één extra vaardigheid: schatten hoeveel keer
de \omterm{def:g5:division:multiple}{deler} past. Schrijf eerst een klein tabelletje met \omterm{def:g5:division:multiple}{veelvouden} van de
\omterm{def:g5:division:multiple}{deler}. Voor $986 \div 23$ (\omterm{def:g5:division:multiple}{veelvouden} van 23: 23, 46, 69, 92, 115, 138, 161,
184, 207):
\begin{enumerate}
  \item $98$ tientallen: $23 \times 4 = 92$ past, $23 \times 5 = 115$ niet:
        cijfer $4$, \omterm{def:g3:division:remainder}{rest} $6$;
  \item haal de $6$ naar beneden: $66$: $23 \times 2 = 46$ past: cijfer
        $2$, \omterm{def:g3:division:remainder}{rest} $20$;
  \item \omterm{def:g3:division:remainder}{quotiënt} $42$, \omterm{def:g3:division:remainder}{rest} $20$. Controle:
        $23 \times 42 + 20 = 966 + 20 = 986$.
\end{enumerate}
\end{method}

\section{Decimale quotiënten}

\begin{example}[De deling die niet bij de rest wil stoppen]\label{ex:g5:division:decimal}
Verdeel $3$ pizza's over $4$ kinderen: $3 \div 4$ heeft \omterm{def:g3:division:remainder}{quotiënt} $0$ en \omterm{def:g3:division:remainder}{rest}
$3$ --- daar heb je niets aan! Ga met de deling verder, voorbij het
\omterm{def:g4:decimals:point}{decimaalteken}:
\begin{enumerate}
  \item $3$ eenheden $= 30$ tienden; $30 \div 4$: $7$ tienden, \omterm{def:g3:division:remainder}{rest} $2$
        tienden;
  \item $2$ tienden $= 20$ honderdsten; $20 \div 4 = 5$ honderdsten, \omterm{def:g3:division:remainder}{rest}
        $0$.
\end{enumerate}
Dus $3 \div 4 = 0.75$: elk kind krijgt $0.75$ pizza --- driekwart, zoals de
taal van de \omterm{def:g4:fractions:def}{breuken} al wist (\cref{prop:g5:fractions:decimal}).
\end{example}

\begin{method}[Een deling voorbij het decimaalteken voortzetten]\label{met:g5:division:continue}
Als de eenheden op zijn en er nog een \omterm{def:g3:division:remainder}{rest} is:
\begin{enumerate}
  \item zet het \omterm{def:g4:decimals:point}{decimaalteken} in het \omterm{def:g3:division:remainder}{quotiënt};
  \item zet een \omterm{def:g1:counting:zero}{nul} achter de \omterm{def:g3:division:remainder}{rest} (eenheden worden tienden, tienden worden
        honderdsten\,\dots) en deel verder;
  \item stop als de \omterm{def:g3:division:remainder}{rest} $0$ is --- of als de cijfers zich eindeloos gaan
        herhalen, zoals bij $1 \div 3 = 0.333\dots$: geef dan een afgeronde
        waarde.
\end{enumerate}
\end{method}

\begin{example}\label{ex:g5:division:continue}
$27 \div 6$: \omterm{def:g3:division:remainder}{quotiënt} $4$, \omterm{def:g3:division:remainder}{rest} $3$; \omterm{def:g4:decimals:point}{decimaalteken}; $30$ tienden
$\div\ 6 = 5$: dus $27 \div 6 = 4.5$. \quad
$22 \div 8$: $2$, \omterm{def:g3:division:remainder}{rest} $6$; dan $60 \div 8 = 7$ \omterm{def:g3:division:remainder}{rest} $4$; dan
$40 \div 8 = 5$: dus $22 \div 8 = 2.75$.
\end{example}

\section{Veelvouden en delers}

\begin{definition}[Veelvoud, deler]\label{def:g5:division:multiple}
De \emph{veelvouden}\index{veelvoud} van $6$ zijn zijn tafel, eindeloos
voortgezet: $0, 6, 12, 18, 24, \dots$ Als $24$ een veelvoud van $6$ is, zeg
je ook dat $6$ een \emph{deler}\index{deler} van $24$ is, of dat $24$
\emph{deelbaar} is door $6$: de deling $24 \div 6$ laat geen \omterm{def:g3:division:remainder}{rest}.
\end{definition}

\begin{proposition}[Deelbaarheidsregels]\label{prop:g5:division:criteria}
Zonder te delen is een getal deelbaar:
\begin{itemize}
  \item door $2$ als zijn laatste cijfer $0$, $2$, $4$, $6$ of $8$ is (de
        \omterm{def:g3:numbers:evenodd}{even} getallen);
  \item door $5$ als zijn laatste cijfer $0$ of $5$ is;
  \item door $10$ als zijn laatste cijfer $0$ is;
  \item door $3$ als de \emph{\omterm{def:g1:addition:def}{som} van zijn cijfers} deelbaar is door $3$
        (bijvoorbeeld $741$: $7 + 4 + 1 = 12$, deelbaar --- en inderdaad
        $741 = 3 \times 247$).
\end{itemize}
\end{proposition}
\admitted

\begin{omfigure}
\begin{tikzpicture}[scale=0.5]
  \foreach \n in {0,...,29} {
    \pgfmathtruncatemacro{\r}{int(\n/10)}
    \pgfmathtruncatemacro{\c}{mod(\n,10)}
    \pgfmathtruncatemacro{\lbl}{\n+1}
    \pgfmathtruncatemacro{\mtwo}{mod(\lbl,2)}
    \pgfmathtruncatemacro{\mthree}{mod(\lbl,3)}
    \ifnum\mtwo=0
      \fill[omDef!25] (\c,-\r) ++(-0.45,-0.45) rectangle ++(0.9,0.9);
    \fi
    \ifnum\mthree=0
      \draw[omThm, thick] (\c,-\r) circle (0.42);
    \fi
    \node[font=\footnotesize] at (\c,-\r) {\lbl};
  }
  \node[font=\small, below] at (4.5,-2.8)
    {blauwe vakjes: veelvouden van $2$;\ \ rode cirkels: veelvouden van
    $3$;\ \ beide: veelvouden van $6$};
\end{tikzpicture}

{\small De getallen $1$ tot $30$: de \omterm{def:g5:division:multiple}{veelvouden} van $2$ en van $3$ vormen
patronen --- en de getallen met beide tekens ($6$, $12$, $18$, $24$, $30$)
zijn precies de \omterm{def:g5:division:multiple}{veelvouden} van $6$.}
\end{omfigure}

\begin{example}\label{ex:g5:division:criteria}
Is $534$ deelbaar door $2$? Het eindigt op $4$: ja. Door $5$? Nee. Door $3$?
$5 + 3 + 4 = 12$: ja. Dus is $534$ ook deelbaar door $6$ (door $2$ \emph{en}
door $3$) --- controle: $534 = 6 \times 89$.
\end{example}

\section{Oefeningen}

\begin{exercise}[$\star$]\label{exo:g5:division:1}
Maak de staartdelingen (eerst een tabelletje met \omterm{def:g5:division:multiple}{veelvouden}):
$851 \div 23$; \quad $704 \div 32$.
\end{exercise}

\begin{exercise}[$\star$]\label{exo:g5:division:2}
Reken de precieze decimale \omterm{def:g3:division:remainder}{quotiënten} uit: $9 \div 2$; \quad $27 \div 4$;
\quad $33 \div 5$; \quad $21 \div 8$.
\end{exercise}

\begin{exercise}[$\star$]\label{exo:g5:division:3}
Vier vrienden verdelen een restaurantrekening van $54$ gelijk. Hoeveel
betaalt ieder? (Ga voorbij het \omterm{def:g4:decimals:point}{decimaalteken} verder.)
\end{exercise}

\begin{exercise}[$\star$]\label{exo:g5:division:4}
Reken $10 \div 3$ uit, en ga drie cijfers voorbij het \omterm{def:g4:decimals:point}{decimaalteken} door.
Wat valt je op? Geef de waarde afgerond op het honderdste.
\end{exercise}

\begin{exercise}[$\star$]\label{exo:g5:division:5}
Noem: de \omterm{def:g5:division:multiple}{veelvouden} van $7$ tot $70$; de \omterm{def:g5:division:multiple}{delers} van $18$; de \omterm{def:g5:division:multiple}{delers} van
$24$.
\end{exercise}

\begin{exercise}[$\star$]\label{exo:g5:division:6}
Deelbaar door $2$? Door $5$? Door $10$? Door $3$? Test elke regel op:
$470$; \quad $735$; \quad $8\,001$; \quad $1\,314$.
\end{exercise}

\begin{exercise}[$\star$]\label{exo:g5:division:7}
Zoek alle getallen tussen $60$ en $80$ die deelbaar zijn door $3$ \emph{en}
door $5$. (Welke enkele deelbaarheid combineert dat?)
\end{exercise}

\begin{exercise}[$\star$]\label{exo:g5:division:8}
Een lint van $7.2$ m wordt in $6$ gelijke stukken geknipt. Hoe lang is elk
stuk? (Deel een decimaal getal door een heel getal: verdeel eerst de meters,
dan de tienden.)
\end{exercise}

\begin{exercise}[$\star$]\label{exo:g5:division:9}
$1\,000$ knikkers worden in zakjes van $24$ verpakt. Hoeveel volle zakjes
krijg je, en hoeveel knikkers blijven over? (Een \omterm{def:g5:division:multiple}{deler} van twee cijfers.)
\end{exercise}

\begin{exercise}[$\star\star$]\label{exo:g5:division:10}
Kunnen $5$ vrienden $7$ chocoladereepjes gelijk verdelen, zonder \omterm{def:g3:division:remainder}{rest}? Geef
elk deel als decimaal getal. Dezelfde vraag voor $7$ vrienden en $5$
reepjes --- wat gaat daar anders?
\end{exercise}

\begin{exercise}[$\star\star$]\label{exo:g5:division:11}
Zoek met de cijfersomregel het kleinste cijfer $d$ waarvoor $52d$ (een getal
van drie cijfers dat op $d$ eindigt) deelbaar is door $3$. Zoek daarna alle
cijfers $d$ die werken.
\end{exercise}
